\section{Preliminaries}
\paragraph{Flow-Matching Models}
Flow Matching (FM) maps a noise distribution $p_0$ to data $p_{\text{data}}$ via an ODE $\text{d}\mathbf{x}_t = v_t(\mathbf{x}_t, t) \text{d}t$. Under the Optimal Transport (OT) formulation, the path is $\mathbf{x}_t = (1-t)\mathbf{x}_0 + t\mathbf{x}_1$, and the model $v_\theta$ learns the constant velocity $(\mathbf{x}_1 - \mathbf{x}_0)$ via:
\begin{equation}
    \mathcal{L}_{\text{FM}}(\theta) = \mathbb{E}_{t, \mathbf{x}_0, \mathbf{x}_1} \left[ \| v_\theta(\mathbf{x}_t, t) - (\mathbf{x}_1 - \mathbf{x}_0) \|^2 \right]
\end{equation}
Following Flow-GRPO~\cite{liu2025flow}, we conceptualize the discretized ODE integration as a sequential \textit{Markovian denoising process}. By formulating each transition $\mathbf{x}_{t} \to \mathbf{x}_{t+\Delta t}$ as a Markovian state step, this perspective bridges continuous generative dynamics with reinforcement learning, defining a formal trajectory for step-wise policy optimization.

\paragraph{On-Policy Distillation}
Knowledge distillation aims to compress teacher capabilities into a student model by minimizing their output divergence. To mitigate distribution shift, on-policy distillation (OPD)~\cite{lu2025onpolicydistillation} requires the student $f_\theta$ to generate trajectories $\tau \sim p_\theta(\tau)$ under the guidance of real-time teacher supervision. For Autoregressive (AR) models, this optimization is formulated as minimizing the Reverse Kullback-Leibler (KL) divergence between the student and teacher distributions:
\begin{equation}
\label{opd}
    \mathcal{L}_{\text{OPD}} = -\mathbb{E}_{y \sim \pi_\theta} \left[ \log \frac{\pi_{\text{teacher}}(y|x)}{\pi_{\theta}(y|x)} \right] = D_{\text{KL}}(\pi_{\theta} \| \pi_{\text{teacher}})
\end{equation}
By aligning the model on its own generated distribution, OPD effectively suppresses exposure bias and ensures robust generalization in interactive or iterative generation tasks.